% from ega0-18-fr.tex
\oldpage[0\textsubscript{IV-1}]{56}
du produit fibré vis-à-vis de chacun des facteurs montre en outre que si l'on
a un morphisme de deux extensions de $B'$
\[
  \xymatrix{
    0\ar[r] & L'_1\ar[r]\ar[d]^h & E'_1\ar[r]\ar[d]^g
      & B'\ar[r]\ar[d]^{1_{B'}} & 0\\
    0\ar[r] & L'_2\ar[r] & E'_2\ar[r] & B'\ar[r] & 0
  }
\]
on en déduit un morphisme \emph{image réciproque par $v$}
\[
  \xymatrix{
    0\ar[r] & L'_1\ar[r]\ar[d]^h & E'_1\times_{B'}B\ar[r]
      \ar[d]^{g\times_{B'}1_B} & B\ar[r]\ar[d]^{1_B} & 0\\
    0\ar[r] & L'_2\ar[r] & E'_2\times_{B'}B\ar[r] & B\ar[r] & 0.
  }
\]
En particulier, si $E'_1$ et $E'_2$ sont des $A$-extensions
$A$-équivalentes de $B'$ par $L'$, leurs images réciproques par $v$ sont des
$A$-extensions $A$-équivalentes de $B$ par $L'$.

La définition du produit fibré montre que lorsque l'on a un morphisme
(18.2.4.1) de $A$-extensions, il se \emph{factorise à travers} l'image
réciproque de $E'$ par $v$~; de façon précise, il existe un
$A$-homomorphisme unique $u_0:E\to F=E'\times_{B'}B$ rendant commutatif le
diagramme
\[
  \xymatrix{
    0\ar[r] & L\ar[r]^j\ar[d]^{w_0} & E\ar[r]^f\ar[d]^{u_0}
      & B\ar[r]\ar[d]^{1_B} & 0\\
    0\ar[r] & L_0\ar[r]\ar[d]^i\ar@{}[d]|{\wr} & F\ar[r]^{p_2}\ar[d]^{p_1}
      & B\ar[r]\ar[d]^v & 0\\
    0\ar[r] & L'\ar[r]_{j'} & E'\ar[r]_{f'} & B'\ar[r] & 0
  }
\]
où $w_0$ est la restriction de $u_0$ à $L$ et
$p_1\circ u_0=u$, $i\circ w_0=w$.
\end{env}

\begin{env}[18.2.6]
\label{0.18.2.6-fr}
Étudions en particulier les images réciproques d'extensions par des
$A$-homomorphismes \emph{surjectifs}. Considérons un $A$-homomorphisme
surjectif $v:B\to B'$, une $A$-extension $F$ de $B$ par un $B$-bimodule $L$,
et l'idéal bilatère $\mathfrak{R}$ de $B$, noyau de $v$ (qui peut être
considéré comme $F$-bimodule au moyen de l'augmentation $F\to B$)~; on a vu
(18.1.7) que tout homomorphisme de $F$-bimodules
$\theta:\mathfrak{R}\to F$ rendant commutatif le diagramme
\[
  \xymatrix{
    & \mathfrak{R}\ar[ld]^\theta\ar[d]\\
    F\ar[r] & B
  }
\]
détermine une extension $E'$ de $B'=B/\mathfrak{R}$ par $L$ dont l'image
réciproque par $v:B\to B'$ est équivalente à $F$, et que toute $A$-extension
$E'$ de $B'$ par $L$ ayant cette dernière propriété s'obtient ainsi (à une
$A$-équivalence près). En outre, pour que deux homomorphismes
$\theta_1,\theta_2$ de $F$-bimodules de $\mathfrak{R}$ dans $F$ donnent deux
  $A$-extensions $A$-équivalentes de $B$ par $L$, il faut et il suffit qu'il
existe une $A$-équivalence $u$ de la $A$-extension $F$ sur elle-même telle que
